\subsection{Starting a new session}
\label{sec:starting-new-session}

A new fitting problem begins with

\begin{verbatim}
/pfit-new <session-name>
\end{verbatim}

This command converts a user description of a scientific fitting problem into a
session specification. The user provides the model equations, data files,
parameter bounds or fixed values, initial conditions, forcing functions, and the
loss definition. The agent's role is to assemble these inputs into the two files
used by the framework, \texttt{inputs/user\_input.yaml} and
\texttt{generated/user\_model.py}.

Before writing those files, the agent must have a complete specification of the
run: the ODE states and their ordering, the trainable and fixed parameters, the
experiments to be fitted, the initial condition and forcing history for each
experiment, the mapping from data columns to observables or inputs, and the
loss used to compare model outputs with measurements. If the user's note is
complete, the session can be created without further questions. If not, the
agent asks only for the missing scientific decisions.

\subsubsection{User problem note}

The user is not expected to write YAML or Python. A concise note with stable
headings is sufficient. For example:

\begin{verbatim}
# Drug response fitting problem

## User model
files = [pulse_low_raw.csv, washout_raw.csv]
states = [D, R]

Initial conditions:
  pulse_low_raw.csv: D = 0, R = 0
  washout_raw.csv:   D = 0, R = 0.3

ODE:
  dD/dt = dose_rate(t) - k_elim D
  dR/dt = k_on D (1 - R) - k_off R

Observables:
  signal_model(t) = baseline + gain R(t)
  recovery_model(t) = R(t)

## Trainable and fixed parameters
k_elim:   0.01  to 10,  log scale
k_on:     0.001 to 10,  log scale
k_off:    0.001 to 10,  log scale
baseline: 50    to 150, linear scale
gain:     1     to 200, linear scale
fixed parameters: none

## Loss function calculation using model outputs and datasets
pulse_low_raw.csv:
  t = time_min
  dose_rate(t) = dose_rate
  signal_obs = signal
  sigma = signal_sd

washout_raw.csv:
  t = time_min
  dose_rate(t) = 0
  recovery_obs = recovery
  sigma = recovery_sd

For each finite measured value:
  r_signal = (signal_model(t_k) - signal_obs,k) / signal_sd,k
  r_recovery = (recovery_model(t_k) - recovery_obs,k) / recovery_sd,k
  loss = sqrt(mean(r^2)) over all finite residuals

Blank measured entries are missing measurements and are excluded.
\end{verbatim}

This form states the model, unknowns, data mapping, and objective in the terms a
scientist naturally uses, while still giving the agent enough information to
construct the internal schema.


\subsubsection{Preparing the data files}

Raw data files may contain different measured quantities. In a multi-experiment
fit, however, the framework expects every fitting CSV to have the same column
layout: the time column first, followed by the same forcing, measurement, and
uncertainty columns in the same order. When an observable is absent from one
experiment, the user keeps the column and leaves its entries blank. The agent
must not modify data files or write derived copies; it identifies the required
layout, and the user creates the framework-ready CSV files.

The following raw files already use consistent names and time units, but they
measure different outputs. The pulse experiment measures \texttt{signal}; the
washout experiment measures \texttt{recovery}. Both outputs enter the loss, so
the canonical layout must contain the union of these measured channels.

\noindent
\begin{minipage}[t]{0.49\textwidth}
\small
\begin{verbatim}
pulse_low_raw.csv

time_min,dose_rate,signal,signal_sd
0,0.0,100.0,4.0
1,1.0,,4.2
2,1.0,113.5,4.4
4,0.0,,4.1
6,0.0,104.5,4.0
\end{verbatim}
\end{minipage}\hfill
\begin{minipage}[t]{0.49\textwidth}
\small
\begin{verbatim}
washout_raw.csv

time_min,dose_rate,recovery,recovery_sd
0,0.0,0.82,0.04
0.5,0.0,0.78,0.04
1.0,0.0,,0.05
2.5,0.0,0.69,0.04
5.0,0.0,0.58,0.03
8.0,0.0,0.51,0.03
\end{verbatim}
\end{minipage}

For fitting, the user prepares both files with the common layout

\begin{verbatim}
time_min,dose_rate,signal,signal_sd,recovery,recovery_sd
\end{verbatim}

The pulse file has no recovery measurements, so the recovery columns are blank.
The washout file has no signal measurements, so the signal columns are blank.
Blank measured values are loaded as \texttt{NaN} and excluded by the masked
loss. The time grids may differ between files, but the time unit must already be
consistent; the framework does not infer or perform unit conversions.

\noindent
\begin{minipage}[t]{0.49\textwidth}
\small
\begin{verbatim}
pulse_low.csv

time_min,dose_rate,signal,signal_sd,recovery,recovery_sd
0,0.0,100.0,4.0,,
1,1.0,,4.2,,
2,1.0,113.5,4.4,,
4,0.0,,4.1,,
6,0.0,104.5,4.0,,
\end{verbatim}
\end{minipage}\hfill
\begin{minipage}[t]{0.49\textwidth}
\small
\begin{verbatim}
washout.csv

time_min,dose_rate,signal,signal_sd,recovery,recovery_sd
0,0.0,,,0.82,0.04
0.5,0.0,,,0.78,0.04
1.0,0.0,,,,0.05
2.5,0.0,,,0.69,0.04
5.0,0.0,,,0.58,0.03
8.0,0.0,,,0.51,0.03
\end{verbatim}
\end{minipage}

Two representative row conversions are:

\noindent
\begin{minipage}[t]{0.49\textwidth}
\small
\begin{verbatim}
pulse_low_raw.csv
raw columns:
time_min,dose_rate,signal,signal_sd
raw row:
2,1.0,113.5,4.4

pulse_low.csv
canonical columns:
time_min,dose_rate,signal,signal_sd,recovery,recovery_sd
canonical row:
2,1.0,113.5,4.4,,
\end{verbatim}
\end{minipage}\hfill
\begin{minipage}[t]{0.49\textwidth}
\small
\begin{verbatim}
washout_raw.csv
raw columns:
time_min,dose_rate,recovery,recovery_sd
raw row:
2.5,0.0,0.69,0.04

washout.csv
canonical columns:
time_min,dose_rate,signal,signal_sd,recovery,recovery_sd
canonical row:
2.5,0.0,,,0.69,0.04
\end{verbatim}
\end{minipage}

This layout lets one loss function be applied to every experiment: residuals
for \texttt{signal} are finite only in the pulse file, residuals for
\texttt{recovery} are finite only in the washout file, and missing channels do
not contribute to the objective.

\subsubsection{Agent review and generated files}

From the note and CSV headers, the agent infers that the two experiments share
one ODE system and one trainable parameter vector, but differ in the initial
condition for \texttt{R} and in their forcing histories. The measured columns are compared to calculated observables: \texttt{signal}
to \texttt{signal\_model = baseline + gain R}, and \texttt{recovery} to
\texttt{recovery\_model = R}. The corresponding uncertainty columns weight the
residuals, and blank measurements require a finite-data mask.

If information is missing, the agent asks targeted questions such as:

\begin{verbatim}
The recovery column appears only in the washout file. Should it measure R(t)?

The signal column is absent from the washout file. Should it be blank there and masked?

Should dose_rate(t) be zero throughout the washout experiment?
\end{verbatim}

Once the specification is complete, \texttt{/pfit-new} writes the session files.
The YAML records the user-prepared CSV files, initial conditions, parameter
bounds, observables, column declarations, and the selected Diffrax integrator.
The agent chooses that integrator from the equations, bounds, and data
timescales, then explains whether the model is being treated as smooth/stiff,
non-smooth/non-stiff, or a conflict that needs model smoothing. The model file
records the ODE right-hand side, forcing interpolation, calculated observables,
and the masked uncertainty-weighted loss over both measured channels. The user then reviews the generated state ordering,
parameter ranges, equations, and data-column usage before continuing to

\begin{verbatim}
/pfit-check <session-name>
\end{verbatim}
